\subsection{Character values with a prescribed support}\label{sec:character-support}
The dimension--trace pair used above is the two-character case of a standard finite-group calculation. Let $\Gamma$ be finite, let $R(\Gamma)$ be its complex representation ring, and write $R(\Gamma)_0$ for its zero adjunction. Character evaluation at $g\in\Gamma$ is a ring homomorphism and extends to
\[
 \widehat{\operatorname{ev}}_g:R(\Gamma)_0\longrightarrow\mathbb C_0,
 \qquad \tau\longmapsto\tau.
\]
Every present representation class maps to a present character value, even when that value is zero.

\begin{proposition}[Characters on a prescribed set of grades]\label{prop:character-support}
Let $E$ be an index set and $B\subseteq E$ a chosen set of present grades. For each $b\in B$, let $f_b$ be a complex class function on $\Gamma$. There is a unique isomorphism class of finite-dimensional genuine representation $V_b$ with character $f_b$ at each present grade if and only if
\[
 m_{b,\chi}=\frac1{|\Gamma|}\sum_{g\in\Gamma}f_b(g)\chi(g^{-1})
 \in\mathbb Z_{\ge0}
 \qquad(\chi\in\operatorname{Irr}(\Gamma)).
\]
The resulting zero-adjoined coefficient is $[V_b]$ on $B$ and $\tau$ on $E\setminus B$. A zero character on a present grade gives the old zero class, not $\tau$. The full family of zero-adjoined character evaluations is injective on $R(\Gamma)_0$.
\end{proposition}
\begin{proof}
The irreducible characters form an orthonormal basis of the class functions, and the displayed inner product computes each multiplicity \citep[\S2.3, Theorem~4 and Corollary~2, p.~16; \S2.5, Theorem~6, p.~19]{SerreCharacters}. Thus nonnegative integral values construct
\[
 V_b=\bigoplus_{\chi\in\operatorname{Irr}(\Gamma)}
 V_\chi^{\oplus m_{b,\chi}},
\]
and character uniqueness proves necessity and uniqueness. All multiplicities vanish exactly for the zero representation. The chosen support convention places this class at a present grade, while absent grades have coefficient $\tau$. Finally, ordinary characters distinguish elements of $R(\Gamma)$, and every extended evaluation distinguishes the new zero from a present class. Together these facts prove injectivity.
\end{proof}

Replacing $\chi(g^{-1})$ by $\chi(g)$ in the formula extracts the multiplicity of the dual irreducible. A constant boundary class function equal to $-2$ is $-2$ times the trivial character and cannot represent a genuine module, since its value at the identity is negative. These are conditional coefficient tests. A prescribed zero grade or a virtual principal part does not supply a new existence theorem for a moonshine module, and the chosen set $B$ cannot be recovered from an ordinary character series after its zero coordinates have been collapsed.
